% ===================================================================
%  TAVOLA N. 10 — Il mare. --- Il porto.
%  Traduction 2026 d'apres texto/io, controlee sur texto/fr et
%  texto/en. Consigne : voir texto/it/10-tavola-01.tex.
%
%  Le francais decale sa numerotation d'un cran a partir de l'alinea 8
%  (il ecrit « 9. --- » sur la cle t10-08-1) : coquille de l'imprimeur
%  francais, conservee dans la transcription. L'italien suit l'ido,
%  qui compte juste.
%
%  LE MEME MOT ARABE PARAIT DEUX FOIS SUR CETTE PLANCHE, SOUS DEUX
%  FORMES ITALIENNES, ET IL NE DESIGNE PAS LA MEME CHOSE.
%
%  L'alinea 10 nomme l'arsenale \textsuperscript{(21)} et
%  l'alinea 16 la darsena \textsuperscript{(53)}, le bassin a
%  flot ou s'abritent les canots. Les deux mots sont le meme :
%  l'arabe « dār aṣ-ṣināʿa », LA MAISON DE LA FABRIQUE, qui nommait le
%  chantier naval des ports musulmans.
%
%  Venise l'a pris deux fois. Une premiere fois par le vénitien
%  « arzanà » — Dante ecrit « l'arzanà de' Viniziani » au chant XXI de
%  l'Enfer — qui a donne darsena, le bassin. Une seconde fois
%  sous la forme savante arsenale, qui est partie dans toute
%  l'Europe et qui designe partout le magasin d'armes, non le bassin.
%
%     UN SEUL MOT ETRANGER, ENTRE DEUX FOIS PAR LE MEME PORT, A DONNE
%     DEUX MOTS QUI NE SE RECONNAISSENT PLUS ET QUI TIENNENT SUR LA
%     MEME PAGE A SIX ALINEAS DE DISTANCE.
%
%  Le tableau 9 venait de montrer un mot qui fait le tour de l'Europe
%  et revient ; celui-ci montre un mot qui entre deux fois par la meme
%  porte et se dedouble. Le tableau 11 allemand avait « Bank », parti
%  meuble et revenu metier. TROIS FIGURES DU MEME PHENOMENE EN TROIS
%  TABLEAUX : un mot n'a pas d'identite, il n'a que des trajets.
%
%  ET TOUT LE RESTE DU PORT PARLE ITALIEN SANS LE DIRE. Le tableau 10
%  allemand avait montre que le vocabulaire allemand et anglais de la
%  mer est bas-allemand, monte de la Hanse. La moitie sud de la
%  navigation europeenne est genoise et venitienne, et elle a fourni :
%     bussola  la boussole, la petite boite
%     molo, cala, banchina, stiva, rotta, cabotaggio
%     pilota, corsaro, fregata, galera, gondola, regata
%     scirocco, libeccio, tramontana, maestrale
%  Le pilota de l'alinea 14 et la fregata de
%  l'alinea 5 sont donc, dans les trente-trois autres colonnes, des
%  mots italiens ; ici ils sont simplement les mots.
%
%  Une seule exception, et elle est du nord : le bordo de la
%  planche est germanique, comme le « Bord » allemand et le « board »
%  anglais. La mer du Nord a donne le flanc du navire, la
%  Mediterranee tout le reste.
% ===================================================================

%%K t10-apar-1 apar
\VUsaut{4.0mm}
\VUtitre{15.0pt}{60}{TAVOLA N. 10}
\VUsaut{3.0mm}
\VUpk{11.4pt}{40}{Il mare. --- Il porto.}
\VUsaut{3.0mm}

%%K t10-01-1 p
1. --- D'estate c'è chi cerca la calma e l'aria fresca in montagna, e chi preferisce
andare al mare. Noi l'anno scorso abbiamo fatto così, perché ci piace molto
l'\VUgras{oceano} \textsuperscript{(1)}.

%%K t10-02-1 p
2. --- Mi piace moltissimo guardare quell'immensa distesa d'acqua che arriva fino
all'\VUgras{orizzonte} \textsuperscript{(2)}, e vedere i grandi \VUgras{piroscafi}
\textsuperscript{(3)} che partono per una lunga \VUgras{traversata} con a bordo i
viaggiatori diretti in America.

%%K t10-03-1 p
3. --- Per questo mio padre, che sa quello che mi piace, per le vacanze sceglie sempre
una spiaggia molto tranquilla vicino a uno dei nostri grandi \VUgras{porti militari}.

%%K t10-03-2 p
Durante il nostro ultimo soggiorno la \VUgras{squadra} è venuta a dare fondo dietro
l'enorme \VUgras{diga foranea} \textsuperscript{(4)} che chiude e protegge il
\VUgras{bacino militare} \textsuperscript{(5)}, e ho potuto ammirare con comodo le
varie \VUgras{navi da guerra}, dal piccolo \VUgras{sommergibile}
\textsuperscript{(10)}, che si immerge e sparisce sott'acqua, all'enorme
\VUgras{corazzata} \textsuperscript{(9)}, che fino a oggi aveva poco da temere se non
gli attacchi della \VUgras{torpediniera} \textsuperscript{(8)}.

%%K t10-tit-2 sub
\VUsaut{3.0mm}
\VUpk{10.2pt}{40}{Le unità piccole.}
\VUsaut{3.0mm}

%%K t10-04-1 p
4. --- Sapeva già trovare con molta \VUgras{abilità} il punto debole della grossa
corazza metallica e piantarci il pericoloso \VUgras{siluro}, la cui esplosione manda a
fondo la nave; ma era ancora meno temibile del sommergibile, perché i cacciatorpediniere
lo raggiungono facilmente.

%%K t10-05-1 p
5. --- Ho potuto vedere anche i veloci \VUgras{incrociatori corazzati}
\textsuperscript{(6)}, quasi invulnerabili come una vera corazzata, diverse
\VUgras{navi trasporto} \textsuperscript{(12)}, tre o quattro \VUgras{cannoniere}
\textsuperscript{(7)}, e infine una \VUgras{fregata} \textsuperscript{(11)}.
\VUgras{La cosa più bella di tutte è vedere quella flotta che leva le ancore.}

%%K t10-06-1 p
6. --- Che differenza di linea fra quelle grandi masse d'acciaio e gli eleganti
\VUgras{tre alberi} o i graziosi \VUgras{velieri} \textsuperscript{(13)} che si usano
ancora nella marina mercantile, e che vedevo entrare in porto a vele spiegate mentre
passeggiavo sul \VUgras{molo} \textsuperscript{(14)}. È vero che perdevano molto del
loro vantaggio quando il \VUgras{vento} era contrario. Dovevano ammainare tutte le vele
per entrare nel bacino, e ci voleva un \VUgras{rimorchiatore} \textsuperscript{(16)} per
andarli a prendere e portarli alla banchina come semplici \VUgras{chiatte}
\textsuperscript{(17)}.

%%K t10-tit-3 sub
\VUsaut{3.0mm}
\VUpk{10.2pt}{40}{Il porto commerciale.}
\VUsaut{3.0mm}

%%K t10-07-1 p
7. --- Quello che rende così interessante il porto di cui parlo è che non ha soltanto
un bacino militare, fatto artificialmente con lavori giganteschi, ma anche uno
splendido \VUgras{porto commerciale} alla \VUgras{foce} di un grande fiume.

%%K t10-08-1 p
8. --- La lingua di terra che separa i due bacini è fortificata e difesa da una linea
di \VUgras{batterie} e di \VUgras{bastioni} che si spingono in mare. Per camminare
sui \VUgras{bastioni} \textsuperscript{(15)} ci vuole un permesso speciale. L'avevo
ottenuto facilmente grazie a mio zio, che è ingegnere ai \VUgras{cantieri navali}
\textsuperscript{(18)}, e sono andato a vedere le navi che si costruiscono sotto
capannoni immensi.

%%K t10-09-1 p
9. --- Ho visto perfino il varo di una corazzata, e mi ricordo del brivido che ha
percorso tutta la folla quando l'enorme massa, lasciata a se stessa, ha cominciato a
scivolare giù per lo scalo ed è venuta a galleggiare sull'acqua del fiume.

%%K t10-10-1 p
10. --- Per arrivare ai cantieri si attraversa la \VUgras{piazza d'armi}
\textsuperscript{(19)}, dove le truppe della guarnigione vengono a esercitarsi ogni
giorno, e si passa su una \VUgras{passerella} \textsuperscript{(20)} di ferro che
collega la piazza d'armi all'\VUgras{arsenale} \textsuperscript{(21)}.

%%K t10-tit-4 sub
\VUsaut{3.0mm}
\VUpk{10.2pt}{40}{L'arsenale e i bacini.}
\VUsaut{3.0mm}

%%K t10-11-1 p
11. --- Con mio zio ho visitato quel grande stabilimento pieno di \VUgras{cannoni}
\textsuperscript{(22)}, di \VUgras{palle} \textsuperscript{(23)} e di
\VUgras{granate}. Ho visto anche la \VUgras{ferriera} \textsuperscript{(24)} dove si
fanno le \VUgras{caldaie} \textsuperscript{(25)} dei piroscafi.

%%K t10-12-1 p
12. --- Lì vicino ci sono i \VUgras{bacini di carenaggio} \textsuperscript{(26)} --- i
bacini a secco --- e i notevolissimi \VUgras{bacini galleggianti}
\textsuperscript{(27)}, dove ho potuto vedere da vicino, su una nave in riparazione,
l'\VUgras{elica} \textsuperscript{(28)}, la \VUgras{chiglia} \textsuperscript{(29)} e
il \VUgras{timone} \textsuperscript{(30)} di un bastimento.

%%K t10-13-1 p
13. --- Il porto commerciale merita di essere studiato quanto quello militare, e ho
passato molte ore a girellare sulle banchine dove sono ormeggiati i \VUgras{vapori a
ruote} \textsuperscript{(31)} che risalgono il fiume, e a guardare la \VUgras{gru a
cavalletto} \textsuperscript{(32)} al lavoro, che solleva carichi enormi come se fossero
piume.

%%K t10-tit-5 sub
\VUsaut{4.0mm}
\VUpk{10.2pt}{40}{Il pilota.}
\VUsaut{3.0mm}

%%K t10-14-1 p
14. --- Ho ammirato i \VUgras{transatlantici} \textsuperscript{(34)} che lasciavano la
rada con le loro \VUgras{bandiere} \textsuperscript{(35)} al vento, guidati da un
abile \VUgras{pilota} \textsuperscript{(36)} che seguiva a meraviglia il
\VUgras{canale} segnato dalle \VUgras{boe} \textsuperscript{(40)} e dai
\VUgras{fanali}, evitando le \VUgras{maone} \textsuperscript{(41)} che con la loro
unica vela quadra si governano con tanta fatica. A \VUgras{prua}
\textsuperscript{(37)} riuscivo a distinguere le \VUgras{cabine}
\textsuperscript{(39)} di seconda classe, e a \VUgras{poppa} \textsuperscript{(38)} i
saloni dei passeggeri di prima.

%%K t10-15-1 p
15. --- Qualche volta ho guardato anche il carico di un \VUgras{piroscafo da carico}
\textsuperscript{(42)}, in cui si ammucchiavano le merci del \VUgras{magazzino}
\textsuperscript{(43)} e della \VUgras{stazione marittima} \textsuperscript{(44)},
portate dai molti treni della \VUgras{ferrovia di banchina}
\textsuperscript{(45)}.

%%K t10-tit-6 sub
\VUsaut{3.0mm}
\VUpk{10.2pt}{40}{Remare per diporto.}
\VUsaut{3.0mm}

%%K t10-16-1 p
16. --- Sono andato spesso in barca nella \VUgras{darsena} \textsuperscript{(53)}, dove
si riparano le piccole \VUgras{barche} \textsuperscript{(46)}, e mi piaceva molto
maneggiare il \VUgras{remo} \textsuperscript{(47)}. Salivo sul \VUgras{ponte}
\textsuperscript{(48)} di una \VUgras{barca a vela} \textsuperscript{(49)},
arrampicandomi per il \VUgras{sartiame} \textsuperscript{(52)} su per l'\VUgras{albero}
\textsuperscript{(51)} fino ai \VUgras{pennoni} \textsuperscript{(50)}, e poi
continuavo la gita in uno \VUgras{skiff} \textsuperscript{(54)} fra le pesanti
\VUgras{barche da pesca} \textsuperscript{(55)}, dove le \VUgras{reti}
\textsuperscript{(56)} stanno ad asciugare.

%%K t10-17-1 p
17. --- Sulla \VUgras{banchina} \textsuperscript{(58)} mi passavano davanti le
\VUgras{merci} \textsuperscript{(59)} più varie, imballate in \VUgras{casse}
\textsuperscript{(60)} o in \VUgras{balle} \textsuperscript{(61)}. Formavano il
\VUgras{carico} \textsuperscript{(63)} delle chiatte, e potenti \VUgras{gru}
\textsuperscript{(62)} le tiravano fuori e le posavano sui \VUgras{camion}
\textsuperscript{(64)}, mentre gli \VUgras{spedizionieri} le dichiaravano e pagavano i
diritti alla \VUgras{dogana} \textsuperscript{(65)}.

%%K t10-18-1 p
18. --- Alla fine, quando arrivavo al \VUgras{ponte girevole} \textsuperscript{(66)} o
alla \VUgras{chiusa} \textsuperscript{(69)}, passando accanto alla \VUgras{draga}
\textsuperscript{(68)} che ripulisce il \VUgras{canale} \textsuperscript{(67)},
sbarcavo sul molo accanto al \VUgras{semaforo} \textsuperscript{(70)}.

%%K t10-19-1 p
19. --- È dall'altra parte di quel molo che accostano le \VUgras{lance}
\textsuperscript{(71)} che portano a terra gli ufficiali della squadra. È bello vedere
i marinai alzare i remi a tempo mentre la barca, portata dall'\VUgras{onda}
\textsuperscript{(72)}, arriva senza fatica alla scaletta. È qui che si osservano
meglio la \VUgras{marea} e il movimento del \VUgras{riflusso} e del \VUgras{flusso}
sulla \VUgras{spiaggia} \textsuperscript{(73)}.

%%K t10-tit-7 sub
\VUsaut{3.0mm}
\VUpk{10.2pt}{40}{Il circolo ufficiali.}
\VUsaut{3.0mm}

%%K t10-20-1 p
20. --- Per tornare a casa ho preso il \VUgras{tram elettrico} \textsuperscript{(74)},
la cui \VUgras{fermata} \textsuperscript{(75)} è al \VUgras{palo} davanti al circolo
ufficiali.

%%K t10-20-2 p
Dal \VUgras{balcone} \textsuperscript{(76)} del circolo, ornato di \VUgras{ancore}
\textsuperscript{(77)}, di cannoni e di palle, ho visto spesso una splendida riunione
di ufficiali di marina: \VUgras{tenenti} \textsuperscript{(78)} con le loro spade,
\VUgras{capitani d'artiglieria} \textsuperscript{(79)} e di \VUgras{fanteria}
\textsuperscript{(80)} con le loro \VUgras{sciabole}. Dietro di loro stava un
\VUgras{marinaio} \textsuperscript{(81)} che portava loro un \VUgras{cannocchiale}
quando volevano vedere che cosa succedeva lontano.

%%K t10-21-1 p
21. --- Ho visto anche giovani \VUgras{guardiamarina} \textsuperscript{(82)} che
prendevano gli ordini dal loro \VUgras{comandante} \textsuperscript{(83)} o
dall'\VUgras{ammiraglio} \textsuperscript{(84)}, mentre un \VUgras{nostromo}
\textsuperscript{(85)} aspettava per portarli a bordo.

%%K t10-22-1 p
22. --- Da quel balcone la vista è magnifica: domina tutta la spiaggia con le sue
\VUgras{tende} \textsuperscript{(86)} e le sue \VUgras{cabine}
\textsuperscript{(87)}, e all'ora del bagno si vedono i \VUgras{bagnanti}
\textsuperscript{(88)} che si divertono allegramente davanti al \VUgras{casinò}
\textsuperscript{(89)}.

%%K t10-23-1 p
23. --- In fondo alla spiaggia la costa forma una piccola \VUgras{penisola}
\textsuperscript{(91)} rocciosa, dove si rompono i \VUgras{frangenti}
\textsuperscript{(92)} e su cui è costruito un \VUgras{faro} \textsuperscript{(93)} che
indica la via ai marinai di notte. Questa penisola è unita alla terraferma solo da un
\VUgras{istmo} \textsuperscript{(90)} molto stretto.

%%K t10-tit-8 sub
\VUsaut{3.0mm}
\VUpk{10.2pt}{40}{Uno sguardo d'insieme sulla costa.}
\VUsaut{3.0mm}

%%K t10-24-1 p
24. --- La \VUgras{scogliera} \textsuperscript{(94)} continua, intaccata dal mare, che
vi ritaglia una serie capricciosa di \VUgras{baie} \textsuperscript{(95)} e di
\VUgras{promontori} (capi), e la rende molto pittoresca.

%%K t10-24-2 p
Sull'\VUgras{altopiano} \textsuperscript{(96)} che domina lo \VUgras{stretto}
\textsuperscript{(98)} ci sono un \VUgras{forte} \textsuperscript{(97)} molto importante
e qualche villa.

%%K t10-25-1 p
25. --- Quello stretto separa dalla terraferma un'\VUgras{isola}
\textsuperscript{(99)} molto pericolosa, circondata di \VUgras{scogli}
\textsuperscript{(100)} e di \VUgras{secche}, dove a volte fanno naufragio barche e
perfino grandi navi. Per quanto pronti siano i soccorsi e per quanto in fretta
arrivino le \VUgras{scialuppe di salvataggio}, questi \VUgras{naufragi} sono terribili,
e nelle \VUgras{burrasche equinoziali} diverse navi sono affondate con tutto
l'equipaggio e il carico.

%%K t10-25-2 p
Malgrado questo vicinato, \VUgras{il porto} è molto frequentato dai marinai.
\VUsaut{6.0mm}
\VUfilet{22mm}
